\section{Related Work}\label{sec:2}

The work reported here sits at the intersection of six research streams: AI-based recommendation systems, semantic matching models, knowledge-driven AI, LLM-based agents, explainable AI, and trustworthy AI.
We briefly review each stream and identify the gap that motivates our system.

\subsection{AI-based recommendation systems}
Modern recommendation systems rely on learned representations, often combined with sparse lexical signals, to produce personalized rankings.
Matrix factorization \citep{koren2009matrix} was the dominant approach in the 2010s; neural collaborative filtering \citep{he2017neural} and deep retrieval \citep{covington2016deep,le2020cold} replaced it in the late 2010s; sequential models \citep{hidasi2015session} added temporal context in the early 2020s.
The dominant production paradigm is a two-tower neural model: one tower encodes the candidate (user, query, resume), the other encodes the job (item, document), and a similarity score is computed in the shared embedding space.
The two-tower paradigm is accurate and scalable but opaque: the similarity score is a single number, and the per-decision factor decomposition is not part of the inference.
In the recruitment domain specifically, recent ESWA-published work \citep{lavi2025careerbert,saito2024leveraging} has shown that the shared-embedding-space approach (CareerBERT, ESCO-aligned) and the multi-behavior preference modeling approach (Saito \& Sugiyama) both reach strong ranking quality on real-world job-corpora, but neither produces a per-decision factor decomposition alongside the score; the present work closes that gap.
A 2025 systematic review of 85 person--job recommendation studies \citep{tang2025explainable} identifies explainability as a primary open challenge and classifies existing methods into data, model, and output layers; the present work contributes to the output layer (component-level explanations) and the data layer (knowledge-grounded retrieval).
Our system retains the two-tower paradigm as one of six channels but adds an explicit composite score that includes non-embedding channels (skill overlap, experience tier, compensation, remote policy), enabling a per-decision factor decomposition that the two-tower paradigm does not provide.

\subsection{Semantic matching models}
Semantic matching is the retrieval-and-ranking core of any system that compares two documents.
Lexical baselines (BM25, TF--IDF) \citep{robertson1995okapi} remain competitive on small document pools because they are sensitive to exact term matches; dense semantic matchers (Sentence-BERT) \citep{reimers2019sentence} generalize across paraphrases but lose the term-match sensitivity; hybrid models fuse both signals and are the strongest in the published benchmarks \citep{cormack2010reciprocal}.
Cross-encoder rerankers are the strongest single model on most benchmarks but are too slow for the hot path of a high-throughput system.
Our system uses BM25 and Sentence-BERT as two of the six channels, fuses them with a learned reciprocal rank fusion, and reports a cross-encoder result in the offline evaluation for completeness (with the cross-encoder disabled by default in the prototype ranking path).

\subsection{Knowledge-driven AI systems}
Knowledge-driven AI systems use a structured representation of domain knowledge (an ontology, a knowledge graph, a vocabulary) to constrain or augment the inference \citep{nikolakopoulos2023knowledge,gentile2014feature}.
In the recruitment domain, the structured representation is typically a skill vocabulary (a normalized set of skill names and aliases), an experience taxonomy (entry / mid / senior / staff), and a job-level classification (full-time / part-time / contract / internship). Standardized competency ontologies such as ESCO \citep{esco2017} and O*NET \citep{onet2024} provide large, curated skill hierarchies, and taxonomy- or ontology-driven skill matching is well established in person--job matching. Our graded matcher is deliberately a \emph{small, hand-built grouping} in this tradition rather than an alignment to a full standard: it is not a new matching mechanism but an auditable, decomposable instantiation whose relation classes and credits are transparent and frozen; we discuss its coarseness and coverage as limitations in \secref{sec:6}.
Our system uses a skill vocabulary as the shared representation across the candidate-side and employer-side components; the vocabulary is consulted by both components and is the only shared state that crosses the role boundary.
The benefit of the vocabulary is that the matching layer can compare two documents on the basis of normalized skill tokens rather than free-text skill descriptions, which reduces the failure mode of two resumes that describe the same skill differently.

\subsection{LLM-based agents}
LLM-based agents decompose complex tasks into tool calls, intermediate reasoning steps, and communication among specialist components \citep{wei2022chain,yao2023react,sumers2024cognitive}.
In the recruitment domain, LLM-based agents are increasingly used for resume parsing (extract structured fields from a free-text resume), job description normalization (extract structured requirements from a free-text posting), and explanation generation (produce a natural-language rationale for a ranking decision).
Recent multi-agent designs for the recruitment domain \citep{lo2025hiring,bhattacharya2025xcui} have demonstrated that LLM-driven agentic frameworks can match or approach human-rater performance on resume screening (Lo et al., 2025, on a hybrid RAG-LLM framework; Bhattacharya \& Verbert, 2025, on the xCUI conversational system with a 20-participant user study).
Our system uses an LLM for parsing and for explanation generation but does not use the LLM as the primary ranker; the LLM is on the cold path (parsing) and on the post-rank path (explanation), not on the hot path (ranking).
The decision reflects a deliberate engineering trade-off: the hot path must be deterministic, fast, and auditable, and current LLM technology does not meet all three requirements at production scale.

\subsection{Explainable AI for recommendation}
Explainable AI for recommendation is a well-developed research stream with both post-hoc and intrinsic approaches \citep{miller2019explanation,gunning2017explainable,tintarev2015explaining,peake2018explanation,chen2022explainable}.
\emph{Post-hoc} approaches train a separate explainer model on the outputs of a black-box ranker; the explainer is not part of the inference and may not be consistent with the ranker.
\emph{Intrinsic} approaches build the explainer into the ranker; the explanation is part of the inference and is consistent with the ranker by construction; foundational work in this direction for recsys includes Tintarev \& Masthoff's design framework for explanation interfaces \citep{tintarev2015explaining} and Peake \& Wang's explanation-mining approach for latent factor models \citep{peake2018explanation}.
The distinction matters in the recruitment domain because intrinsic approaches provide the engineering prerequisite for an auditable explanation without further engineering, whereas post-hoc approaches must be re-validated whenever the underlying ranker is updated.
In the recruitment domain specifically, Wanner et al. \citep{wanner2020generating} demonstrate that counterfactual explanations can be generated for explainable AI in recruitment, and L\"ofstr\"om et al. \citep{lofstrom2024calibrated} introduce a calibrated-explanations framework that adds uncertainty information to rule-based explanations; the present work combines both directions in a single system and extends the calibration to a multi-channel composite rather than a single classifier.
Our system uses an intrinsic approach: the rule-based explainer is bound to the same six channels that produce the score, and the explanation bullet that names a channel is guaranteed to correspond to that channel's contribution to the score.
Counterfactual explanations \citep{kaur2020interpreting,wanner2020generating} are a complementary approach that explains a ranking by identifying the smallest change to the input that would change the output; our system includes a counterfactual probe that validates whether the predicted change is realized in the next ranking.
The combination of an intrinsic per-decision factor decomposition, a calibrated confidence display, and a counterfactual probe is, to our knowledge, novel in the recruitment-recommendation literature.

\subsection{Trustworthy AI and confidence calibration}
Trustworthy AI is the research stream that treats trust as a property of the inference, not of the model's reputation \citep{liu2020trust,pu2007trust}.
A trustworthy ranking system reports, alongside the ranking, a calibrated confidence value that the user can read as a signal of the system's uncertainty.
The standard measure of calibration quality is the expected calibration error (ECE) \citep{guo2017calibration}, and the standard recalibration method is Platt scaling \citep{platt1999probabilistic} or temperature scaling.
The recent calibrated-explanations framework of L\"ofstr\"om et al. \citep{lofstrom2024calibrated} provides a principled approach to adding uncertainty information to rule-based explanations; we adopt the same conceptual approach of pairing a calibrated confidence with a factor decomposition, applying standard scalar-score calibration to the composite output.
Our system applies Platt scaling to the composite ranking output and reports the calibration parameters (slope and intercept) and the Brier score \citep{brier1950verification} alongside the ECE.
We apply standard Platt scaling to the scalar composite score; because the calibrator operates on the single final score, this is ordinary ranker-score calibration rather than a new calibration mechanism, and we make no novelty claim for the act of calibrating a multi-channel score. Our calibration contribution is instead the honest calibration-versus-discrimination trade-off analysis reported in \secref{sec:5} (Platt attains a low equal-width ECE but near-base-rate discrimination, while beta calibration attains the lowest ECE \emph{and} preserves discrimination; see \secref{sec:5.3}).
Fairness in information access \citep{ekstrand2021rethinking} provides a complementary lens: ranking systems in employment contexts can encode proxy variables that disproportionately disadvantage protected groups, and the engineering remedy is to test for disparate impact under controlled perturbations of the input.
Causal and counterfactual debiasing is a complementary direction for reducing disparate impact in recommendation; our system reports an adversarial fairness probe as an engineering sanity check and a foundation for the larger corpus study described in \secref{sec:7}.
The HCI literature on explainable AI user experiences \citep{amershi2019guidelines,liao2020questioning} provides design guidance for the user-facing explanation interface; the present system follows the same engineering pattern of surfacing a calibrated confidence value and a per-decision factor decomposition in the primary user interface, and reports the user-facing surface in \secref{sec:1.7}.

\subsection{Positioning}
The present work is positioned as an applied AI system for the recruitment domain that integrates six research streams into a single, evaluated pipeline.
The closest prior systems address one or two of the streams but not all six; the comparison is taken up in \secref{sec:6}.
The contribution is the integration, the engineering surface, and the controlled evaluation; the individual components (BM25, Sentence-BERT, RRF, Platt scaling, rule-based explanation) are well-known, and the novelty is in their composition and in the evaluation methodology that maps each component to a measurable property.
